% COMP9517 26T2 Group Project Report — iNaturalist-2021 fine-grained classification
% REQUIRED: official CVPR LaTeX template. Drop cvpr.sty + supporting files into this
% folder (download from the CVPR template linked in the course spec), then compile main.tex.
% Rubric-mapped sections below. Main text <= 10 pages incl. figures/tables (refs excluded).
\documentclass[10pt,twocolumn,letterpaper]{article}
\usepackage[pagenumbers]{cvpr}      % <-- official CVPR template, page numbers on
\usepackage{graphicx,amsmath,amssymb,booktabs,subcaption,xcolor}
\usepackage[pagebackref,breaklinks,colorlinks]{hyperref}

\def\cvprPaperID{****}
\begin{document}

\title{Fine-Grained Species Classification on iNaturalist-2021:\\
Traditional and Deep Pipelines with a Robustness and Explainability Study}
\author{Group NN \quad COMP9517 Computer Vision, 2026 Term 2\\
{\tt\small z5633314, \dots}}
\maketitle

\begin{abstract}
% 4--6 sentences: task, the two methods, headline numbers, the two advanced studies,
% one-line finding. Fill LAST once numbers exist.
We study fine-grained visual classification on a 500-class subset of iNaturalist-2021,
comparing a traditional Bag-of-Visual-Words (SIFT)+SVM pipeline against convolutional
networks trained from scratch and from ImageNet-pretrained weights. \dots
\end{abstract}

%--------------------------------------------------------------------------------
\section{Introduction}\label{sec:intro}
% [rubric 1pt] relevance of automated species recognition; understanding of the task
% (fine-grained, high intra-class variation, long tail) and the dataset. State contributions.

\section{Literature Review}\label{sec:lit}
% [rubric 2pt] handcrafted features (SIFT/HOG/LBP, BoVW), classical classifiers (SVM/RF);
% CNNs + transfer learning; fine-grained techniques; Grad-CAM \cite{gradcam}; corruption
% robustness \cite{hendrycks}. Cite everything you use.

\section{Methods}\label{sec:methods}
% [rubric 4pt] motivate BOTH methods.
\subsection{Traditional pipeline: BoVW-SIFT + SVM}
% descriptor -> KMeans codebook (k words) -> histogram -> LinearSVC. Design choices.
\subsection{Deep pipeline: CNN (scratch vs.\ pretrained)}
% architecture, from-scratch vs ImageNet-pretrained fine-tuning, augmentation, optimiser.

\section{Experimental Results}\label{sec:exp}
% [rubric 4pt] setup: exact subset (500 classes, 40/10/10, seed=42), hyperparameters.
\subsection{Setup and metrics}
% top-1/top-5, overall + balanced acc, macro P/R/F1, train/test time. Emphasise macro-F1.
\subsection{Main comparison}
% Table: method x {top1, top5, macroF1, train_s, test_s}.
\subsection{Ablations}
% one factor at a time: scratch vs pretrained; architecture; augmentation on/off; codebook k.
% training curves (loss/acc vs epoch); confusion matrix; hardest confused species pairs.

\subsection{Advanced study I: Test-time degradation robustness}\label{sec:robust}
% model fixed, degrade ONLY test set. 5 corruptions x several severities.
% Figure: top-1 & macro-F1 vs severity per corruption. Which method is most robust and why.

\subsection{Advanced study II: Grad-CAM explainability}\label{sec:cam}
% organism vs background; correct vs wrong; confusable pairs. Interpret, don't just show maps.

\section{Discussion}\label{sec:disc}
% [rubric 1pt] compare methods; why the winner wins; failure cases with reasons;
% robustness vs image representation; what Grad-CAM revealed.

\section{Conclusion}\label{sec:conc}
% [rubric 1pt] what worked / didn't, limitations, future work.

{\small\bibliographystyle{ieeenat_fullname}\bibliography{refs}}
\end{document}
